\vpara{Coarse-Grained SP-based Methods.}
Many efforts are paid to solve generalization problems on SP-based KBQA. Some approaches, such as~\cite{lan2020query,gu2021beyond}, use a rank-based model that takes advantage of a coarse-level match between the question and the logical expressions or query graphs. They first enumerate numerous query graph candidates based on KBs and then they rank them according to how relevant they are to the question. Another line of approaches, in addition to the rank-based ones, makes use of a generation model. KQAPro~\cite{cao2022kqa} leverages BART to directly convert questions into logical expressions. Additionally, RNG-KBQA~\cite{ye2022rng} further injects top-k ranked logical expressions as an additional input to the question. CBR-KBQA~\cite{das2021case} injects analogous questions and their corresponding logical expressions from the training data to increase the generalization. All of the aforementioned methods are pure coarse-level frameworks that treat each coarse-grained logical expression as a separate unit. 

\vpara{Fine-Grained SP-based Methods.}  
Many researchers have been motivated to address the generalization issue by the notion of utilizing decomposed components, such as class, relation, and logical skeleton.
Some approaches~\cite{wang-etal-2020-rat,zhao-etal-2022-bridging,li2022resdsql} retrieve the relevant schema item such as relation and column as additional fine-grained input information, while another line of approaches~\cite{dong-lapata-2018-coarse} extracts the skeleton of logical expression as the decoder guide. Such methods primarily concentrate on the grammar of logical expression and often ignore the knowledge constraint, which is essential in large-scale KB. They usually focus on KBs or DBs that contain a small number of relations where a logical expression can be easy to be executable. Program Transfer~\cite{cao-etal-2022-program}, ReTrack~\cite{chen-etal-2021-retrack}, and TIARA~\cite{shu2022tiara} simply apply KB constraints to control the generation of the decoding process. As opposed to them, we make use of middle-grained KB constraints during both the encoding and the decoding processes to help the model better adapt to KB and ensure executability. 

